\section{初始实验与公开实例}
\label{sec:03_05}

初始实验用于检查前述模型和求解方法在不同输入结构上的表现。输入既包括可以逐项核验的教学算例，也包括具有配对关系的合成条件和公开论文使用的完整拓扑。解释这些实验之前，需要先明确任务、资源、收益和输入顺序如何形成，以及不同条件之间共享了哪些信息。相同拓扑采用不同权重可以形成不同优化问题，但不能因此计为两个独立数据集；使用公开输入也不等于复现了提供该输入的论文算法。

\subsection{条件组成与统一规模口径}

初始实验共包含$33$个条件，其组成是$1$个十二任务教学算例、$18$个初始合成条件、$6$个公开最大独立集条件、$7$个JGAA公开拓扑及其原档案权重条件，以及$1$个采用ESA代码规则重建权重的条件。最后一个条件与前面的一个JGAA条件共用\texttt{superblue19}拓扑。因此，$33$是本节的实验条件数，不是已经确认的独立数据集数。

对每个条件，资源集合记为$V$，候选任务索引集合记为$E$。每项任务$e\in E$占用非空资源集合$R_e\subseteq V$，选取变量$x_e\in\{0,1\}$表示是否接受该任务。不同任务即使具有相同资源集合，仍保留各自的任务索引，不因集合相同而自动合并。全部初始条件采用资源容量$b_v=1$，$v\in V$，并求解
\begin{equation}
 \begin{aligned}
 \max_{x}\quad &\sum_{e\in E}w_ex_e,\\
 \text{s.t.}\quad &\sum_{e\in E}\mathbb{I}_{\{v\in R_e\}}x_e\leq1,
 &&\forall v\in V,\\
 &x_e\in\{0,1\},&&\forall e\in E.
 \end{aligned}
 \label{s0305:common_model}
\end{equation}
其中$w_e>0$是固定的无量纲收益，指标函数$\mathbb{I}_{\{v\in R_e\}}$在任务$e$使用资源$v$时取$1$，否则取$0$。资源容量和占用量均为不可分资源单位的计数；这里的收益不代表货币、能量或实际巡检质量。

为保持与已有实验记录一致，任务数使用$n$，资源数使用$m$，不随公开文件的头部命名改变。定义
\begin{align}
 n&=|E|, &m&=|V|,\label{s0305:sizes}\\
 r&=\max_{e\in E}|R_e|,
 &\Delta&=\max_{v\in V}\sum_{e\in E}\mathbb{I}_{\{v\in R_e\}}.
 \label{s0305:rank_degree}
\end{align}
四个量均为无量纲整数。$r$表示单项任务最多需要多少种资源，$\Delta$表示竞争同一种资源的任务最多有多少项。前者描述任务的资源需求规模，后者描述容量约束关联的变量规模，二者不能互相替代。

表\ref{tab:initial_parameters}逐行列出全部初始条件。表中参数来自原输入核验，公开实例保持完整拓扑，没有根据求解结果截取局部子图。

{\small\setlength{\tabcolsep}{4pt}
\begin{longtable}{llrrrr}
\caption{初始条件的对象规模和关联结构}\label{tab:initial_parameters}\\
\toprule
编号 & 条件 & $n$ & $m$ & $r$ & $\Delta$\\
\midrule
\endfirsthead
\toprule
编号 & 条件 & $n$ & $m$ & $r$ & $\Delta$\\
\midrule
\endhead
\midrule
\multicolumn{6}{r}{续下页}\\
\endfoot
\bottomrule
\endlastfoot
C01 & 十二任务 & 12 & 12 & 6 & 4\\
C02 & 合成，种子0 & 300 & 200 & 6 & 20\\
C03 & 合成，种子0 & 1000 & 200 & 6 & 46\\
C04 & 合成，种子0 & 3000 & 200 & 6 & 115\\
C05 & 合成，种子0 & 1000 & 200 & 3 & 27\\
C06 & 合成，种子0 & 1000 & 200 & 12 & 81\\
C07 & 合成，种子0 & 1000 & 200 & 24 & 145\\
C08 & 合成，种子1 & 300 & 200 & 6 & 17\\
C09 & 合成，种子1 & 1000 & 200 & 6 & 44\\
C10 & 合成，种子1 & 3000 & 200 & 6 & 116\\
C11 & 合成，种子1 & 1000 & 200 & 3 & 26\\
C12 & 合成，种子1 & 1000 & 200 & 12 & 83\\
C13 & 合成，种子1 & 1000 & 200 & 24 & 146\\
C14 & 合成，种子2 & 300 & 200 & 6 & 19\\
C15 & 合成，种子2 & 1000 & 200 & 6 & 49\\
C16 & 合成，种子2 & 3000 & 200 & 6 & 117\\
C17 & 合成，种子2 & 1000 & 200 & 3 & 25\\
C18 & 合成，种子2 & 1000 & 200 & 12 & 80\\
C19 & 合成，种子2 & 1000 & 200 & 24 & 151\\
C20 & 公开MIS，度3 & 500 & 750 & 3 & 2\\
C21 & 公开MIS，度3 & 1000 & 1500 & 3 & 2\\
C22 & 公开MIS，度3 & 2000 & 3000 & 3 & 2\\
C23 & 公开MIS，度5 & 500 & 1250 & 5 & 2\\
C24 & 公开MIS，度5 & 1000 & 2500 & 5 & 2\\
C25 & 公开MIS，度5 & 2000 & 5000 & 5 & 2\\
C26 & ibm01 & 14111 & 12752 & 42 & 39\\
C27 & ibm02 & 19584 & 19601 & 134 & 69\\
C28 & us04 & 163 & 28016 & 11625 & 18\\
C29 & bibd\_49\_3 & 1176 & 18424 & 47 & 3\\
C30 & AProVE07-01 & 7502 & 28770 & 848 & 847\\
C31 & AProVE07-27 & 7729 & 29194 & 3246 & 201\\
C32 & superblue19 & 511685 & 522482 & 9155 & 1507\\
C33 & superblue19（LHG） & 511685 & 522482 & 9155 & 1507\\
\end{longtable}
}
相同公开拓扑的不同权重条件分别列出，不按不同数据集计数。


教学条件C01连接第2章的数值模型与本章的求解实现，其作用是允许读者追踪任务选择和资源占用。其余条件扩展了变量规模与关联结构，但不改变式\eqref{s0305:common_model}中的单位容量假设。

\subsection{合成条件的前缀配对与嵌套资源需求}

C02至C19均使用固定的$200$种资源。每个种子$s\in\{0,1,2\}$对应以下六组任务数与秩设置：
\begin{equation}
 (n,r)\in\{(300,6),(1000,6),(3000,6),
 (1000,3),(1000,12),(1000,24)\}.
 \label{s0305:synthetic_settings}
\end{equation}
每个设置保留三个预先固定的种子，共形成$18$个条件。这些种子用于生成输入，不代表在一个固定输入上重复求解三次。

构造过程使用彼此分离的资源抽样流和收益抽样流。资源流的伪随机种子为$814000+s$，收益流的伪随机种子为$927000+s$。令资源编号集合为$V=\{0,\ldots,199\}$。对任务编号$i\in\{0,\ldots,2999\}$，先从$V$中无放回抽取一个含$24$个不同资源的有序列表
\begin{equation}
 B_i^{(s)}=(v_{i,1}^{(s)},\ldots,v_{i,24}^{(s)}),
 \qquad v_{i,j}^{(s)}\in V,
 \label{s0305:resource_list}
\end{equation}
其中同一列表内不同位置对应不同资源；同时生成整数收益$w_i^{(s)}\in\{1,\ldots,100\}$。对于给定的$(n,r)$，只使用前$n$个任务，并令
\begin{equation}
 E_n=\{0,\ldots,n-1\},\qquad
 R_i^{(s,r)}=\{v_{i,1}^{(s)},\ldots,v_{i,r}^{(s)}\},
 \quad i\in E_n.
 \label{s0305:synthetic_input}
\end{equation}
资源列表在写入时按编号排序，但任务流保持任务编号顺序；配对方法接收相同的任务顺序和收益。抽样流分离使资源需求长度的改变不会推进收益抽样序列，从而避免同时改变同一任务的收益。

该构造提供两种明确的配对关系。第一，在$r=6$时，$300$任务条件和$1000$任务条件分别是$3000$任务条件的完整前缀，已有任务的资源需求和收益均保持不变。第二，在$n=1000$时，同一任务满足
\begin{equation}
 R_i^{(s,3)}\subseteq R_i^{(s,6)}
 \subseteq R_i^{(s,12)}\subseteq R_i^{(s,24)},
 \qquad i\in E_{1000},\ s\in\{0,1,2\},
 \label{s0305:nested_resources}
\end{equation}
并在四个条件中使用相同收益$w_i^{(s)}$。因此，增加$r$具有逐项增加既有任务资源需求的含义，而不是重新生成一组不相关任务。

这些配对关系也限定了实验的解释。对同一$s$和$r=6$，较小输入的任意可行解都可以通过令新增任务变量为零，扩展为较大输入的可行解。若$z_{n,r,s}^{*}$表示相应条件的最优收益，则
\begin{equation}
 z_{300,6,s}^{*}\leq z_{1000,6,s}^{*}
 \leq z_{3000,6,s}^{*},\qquad s\in\{0,1,2\}.
 \label{s0305:size_optimum_relation}
\end{equation}
另一方面，在固定任务数和收益时，式\eqref{s0305:nested_resources}只会保留或增加冲突。因此，记$\mathcal{F}_{n,r,s}\subseteq\{0,1\}^{n}$为对应可行域，有
\begin{align}
 \mathcal{F}_{1000,24,s}&\subseteq\mathcal{F}_{1000,12,s}
 \subseteq\mathcal{F}_{1000,6,s}\subseteq\mathcal{F}_{1000,3,s},
 \label{s0305:rank_feasible_relation}\\
 z_{1000,24,s}^{*}&\leq z_{1000,12,s}^{*}
 \leq z_{1000,6,s}^{*}\leq z_{1000,3,s}^{*}.
 \label{s0305:rank_optimum_relation}
\end{align}
上述关系由输入构造直接推出，是最优值和可行域的数学关系，不是对限时算法输出的单调性承诺。

在这些合成条件中，每项任务恰好关联$r$种资源，故总关联项数和每种资源的平均关联任务数分别为
\begin{equation}
 \sum_{v\in V}\sum_{i\in E_n}
 \mathbb{I}_{\{v\in R_i^{(s,r)}\}}=nr,
 \qquad
 \frac{1}{200}\sum_{v\in V}\sum_{i\in E_n}
 \mathbb{I}_{\{v\in R_i^{(s,r)}\}}=\frac{nr}{200}.
 \label{s0305:synthetic_incidence}
\end{equation}
因此，在资源数固定时增加$n$，不仅增加决策变量，也增加平均资源竞争；增加$r$则同时改变关联项数、冲突结构和可行域。求解时间还受预处理、启发式选择和最优性证明过程影响，不能仅由$n$或$r$增加就推断计算难度必然单调增加。三个固定种子的输入用于开发和配对诊断，也不足以支持总体显著性结论。

\subsection{公开MIS输入的等价转换}

C20至C25来自Li等再利用研究关联的六个完整公开最大独立集输入\cite{LiHypOp2025}。原图顶点数为$500$、$1000$和$2000$，原图度数$d$分别为$3$和$5$。这里$d$表示原图中每个顶点的邻接边数，不是神经网络宽度或其他学习参数。输入核验包括图头部计数、边数、端点范围、无自环和重复边、各顶点度数，以及原作者日志中的文件对应关系。

为避免与转换后的资源集合混淆，原图记为$G=(Q,F)$，其中$Q$是原图顶点集合，$F$是无向边集合。对每个原顶点$q\in Q$建立候选任务$e_q$；对每条原边$f=\{p,q\}\in F$建立容量为$1$的资源$v_f$。转换定义为
\begin{align}
 E&=\{e_q:q\in Q\},&
 V&=\{v_f:f\in F\},\label{s0305:mis_objects}\\
 R_{e_q}&=\{v_f:f\in F,\ q\in f\},&
 w_{e_q}&=1,\qquad q\in Q.
 \label{s0305:mis_mapping}
\end{align}
原图一条边的两个端点任务恰好共享该边对应的资源。因而式\eqref{s0305:common_model}中的资源约束逐条成为
\begin{equation}
 x_{e_p}+x_{e_q}\leq1,
 \qquad \forall f=\{p,q\}\in F.
 \label{s0305:mis_constraints}
\end{equation}
每个约束禁止同时选择原图中的两个相邻顶点，没有增加原图中不存在的邻接限制。

具体地，设$S\subseteq Q$是原图独立集，并令$M(S)=\{e_q:q\in S\}$。由于$S$不含任何一条原边的两个端点，$M(S)$在每个资源上至多占用一次，故是可行任务集。反过来，若$M\subseteq E$满足所有容量约束，则$S(M)=\{q\in Q:e_q\in M\}$不能包含相邻顶点，否则相应资源会被占用两次。因此$S(M)$是独立集。两个映射互为逆映射，构成独立集与可行任务集之间的双射，并且
\begin{equation}
 \sum_{e\in M(S)}w_e=|M(S)|=|S|,
 \qquad \forall S\subseteq Q\text{ 为独立集}.
 \label{s0305:mis_objective}
\end{equation}
由此可知，转换既保持可行解，也保持每个解的收益与最优值；并非只保持某一个已知解。

这些输入是$d$正则简单图，故转换后
\begin{equation}
 n=|Q|,\qquad m=|F|=\frac{nd}{2},\qquad
 r=d,\qquad\Delta=2.
 \label{s0305:mis_sizes}
\end{equation}
等式$m=nd/2$来自每条原边具有两个端点。$r=d$说明每个任务需要原顶点关联的$d$种资源，而$\Delta=2$说明每种资源始终只连接两个任务。这里较大的任务秩并未产生超出原图两两冲突的约束。

六个图均完整保留，任务按原顶点编号升序输入，没有重新抽样或截取。实验复用了HypOp相关公开输入，并在等价模型上比较SCIP与Streaming；HypOp神经网络求解器未运行。因此，这部分结果只能解释为公开输入上的新求解比较，不能称为HypOp算法复现，也不能据此比较神经求解器与本章方法的性能。

\subsection{JGAA完整拓扑与单位容量条件}

C26至C32对应Großmann等研究关联的七个完整公开拓扑\cite{GrossmannJRS2026}，分别为\texttt{ibm01}、\texttt{ibm02}、\texttt{us04}、\texttt{bibd\_49\_3}、\texttt{AProVE07-01}、\texttt{AProVE07-27}和\texttt{superblue19}。这些输入保留档案中的全部超边、资源成员及其顺序，超边权重采用原档案中$1$至$100$的随机整数权重，没有重新生成权重或按求解结果重排任务。

公开超图文件同时保存了超边权重和顶点权重。本节将超边解释为候选任务，将超图顶点解释为资源；档案中的随机顶点权重保留在源文件中，但不作为本次资源容量。当前实验明确选择论文中的统一单位容量条件，即
\begin{equation}
 b_v=1,\qquad \forall v\in V.
 \label{s0305:public_capacity}
\end{equation}
因此，本节输入保留原拓扑和原超边收益，同时明确指定容量为$1$。它不是对档案所有字段含义不加区分的直接求解，也不覆盖论文中的其他容量条件。

公开条件的规模必须按表\ref{tab:initial_parameters}中的同一行理解。例如，\texttt{us04}具有$n=163$和$r=11625$，说明任务数不多，但个别任务涉及大量资源；\texttt{superblue19}具有$n=511685$、$m=522482$、$r=9155$和$\Delta=1507$，同时具有较大的任务数和资源数。不能把前者的最大秩与后者的任务数拼接成一个不存在的实例。最大秩还只描述最宽的一项任务，不能据此认为全部任务均具有该资源需求规模。

这些条件的选择面向规模和结构的代表性观察，并非对原论文全部数据和配置的复现。第\ref{sec:03_04}节的原作者完整流程也使用这七个已有条件，其$10$个固定种子是算法运行重复，不是新增输入拓扑。相应$70$次计划运行中保留了$65$次有效运行和$5$次失败，失败不能因输入属于公开基准而省略，也不能通过只报告有效子集来改变计划运行数。

\subsection{ESA共用拓扑与权重重建}

C33与C32使用相同的完整\texttt{superblue19}拓扑，但采用Reinstädtler等半流式匹配研究公开代码中的权重构造规则\cite{ReinstadtlerFPU2025}。设该拓扑的最大任务资源数为$r_{\max}=\max_{e\in E}|R_e|$，则重建权重为
\begin{equation}
 w_e^{\mathrm{ESA}}=r_{\max}+1-|R_e|,
 \qquad \forall e\in E.
 \label{s0305:esa_weights}
\end{equation}
$r_{\max}$与本节的$r$具有相同计数含义；保留下标是为了明确权重由该完整拓扑的最大秩计算。由于$1\leq|R_e|\leq r_{\max}$，重建权重为正整数，且资源需求最大的任务仍具有权重$1$。该规则在同一拓扑内赋予资源需求较少的任务较高收益。

这里需要保留论文文字、公开代码与可获得文件之间的差异。论文文字给出的表达省略了式\eqref{s0305:esa_weights}中的$+1$，本实验依据公开代码重建权重；原始生成的ESA加权文件没有取得。因此，C33应表述为完整公开拓扑上的公开代码权重重建条件，不能称为已经找到并直接使用的完全原始ESA加权数据。

这个差异也不是可以无条件忽略的整体目标平移。若将省略$+1$的文字表达记为$\widetilde w_e=r_{\max}-|R_e|$，则对任意可行任务集$M\subseteq E$有
\begin{equation}
 \sum_{e\in M}w_e^{\mathrm{ESA}}
 =\sum_{e\in M}\widetilde w_e+|M|.
 \label{s0305:esa_shift}
\end{equation}
不同可行解的任务数$|M|$可以不同，故逐项增加$1$可能改变解之间的收益排序；省略$+1$还会使最大秩任务的权重变为零。采用代码规则是一个必须公开说明的输入定义，不能据此宣称两种权重规则具有相同最优解。

C32与C33的资源集合、任务关联和容量完全一致，因而可行域相同，目标系数则不同。两行分别报告可以观察收益构造对求解行为的影响，但它们共享拓扑，不能视为两个相互独立的结构样本。

\subsection{初始求解口径与结论边界}

初始比较采用SCIP~10.0.2及PySCIPOpt~6.2.1建立单位容量匹配整数规划，SCIP的初始求解预算为每条件$10\,\mathrm{s}$、单线程；Streaming采用固定的$\delta_g$配置，$\varepsilon=0.1$，按给定顺序单遍处理并执行原有逆栈输出。Streaming在全部$33$个初始条件上运行。$10\,\mathrm{s}$是SCIP初始阶段的预算上限，并非两种方法都必须消耗的固定时间，也不包含任意后续继续求解的成本。

历史记录还包含部分条件的首次正收益解，以及后续按$5\%$间隙条件停止的记录。这些记录属于不同求解阶段。首次正解只表明已找到正收益可行候选；达到给定间隙停止条件也不等于证明全局最优。第\ref{sec:03_07}节的结果表列示已完成记录时，需要随值保留相应阶段和停止口径，不能把后续阶段的收益配上初始$10\,\mathrm{s}$预算，形成并不存在的成本与质量组合。求解器运行中的数值快照也不能替代最终任务编号、容量和收益的独立验收。

在最大化问题中，设某一已核验可行解收益为$L$，该次求解提供的有限有效上界为$U$，相应最优收益为$z^*$，则
\begin{equation}
 L\leq z^*\leq U.
 \label{s0305:bound_meaning}
\end{equation}
$L$、$U$与$z^*$均使用该条件的无量纲收益尺度。限时停止而未证明最优时，其他方法相对于$L$的收益差只是与当前可行解的差异，不能称为相对于最优值的损失。若缺少有限有效上界，则不能仅由当前可行收益补造最优性证书。

初始协议将建模或转换、求解以及结果提取纳入端到端计算口径，并另列求解阶段时间；磁盘读写、导入和外部核验不计入这一口径。后续完整流程和其他历史阶段的停止条件、热启动及CPU并发条件存在差别，因此不同批次的时间只能按各自定义解释，不能直接合并成等预算速度排名。

本节输入服务于模型开发、配对诊断和代表性公开实例观察。合成条件共享前缀或嵌套资源需求，公开条件经过代表性选择，且存在共用拓扑的权重变体；这些特点决定了它们不能被称为正式独立泛化测试。此外，全部条件均满足单位容量，任意共享资源的两项任务互斥，其二元可行域仍可由冲突图表示。较大的$r$能够说明单项任务的资源关联规模，却不能单独证明不可约的高阶依赖或超图算法的性能优势。
